\documentclass{beamer}

%%%%%%%%%%%%%%%%%%%%%%%%%%%%%
% CONFIGURACIÓN DEL TEMA Y COLORES
%%%%%%%%%%%%%%%%%%%%%%%%%%%%%
\usetheme{Madrid}            % Tema principal
\usecolortheme{seahorse}     % Combinación de colores
\setbeamertemplate{navigation symbols}{} % Quitar la barra de navegación
\setbeamertemplate{footline}[frame number] % Mostrar el número de diapositiva
\setbeamertemplate{itemize items}[circle]  % Puntos de las listas en forma de círculo

%%%%%%%%%%%%%%%%%%%%%%%%%%%%%
% DATOS DE LA PRESENTACIÓN
%%%%%%%%%%%%%%%%%%%%%%%%%%%%%
\title{\textbf{Introducción y Marco Legal en PRL \\ \Large(Prevención de Riesgos Laborales)}}
\subtitle{Incluyendo medidas de seguridad y protección medioambiental}
\author{\textbf{Instructor:} Santo Domingo Savio}
\date{\today}
\usepackage{pgfpages}
%\setbeameroption{show notes on second screen=left}

%%%%%%%%%%%%%%%%%%%%%%%%%%%%%
% DOCUMENTO
%%%%%%%%%%%%%%%%%%%%%%%%%%%%%
\begin{document}

%----------------------------------
% PORTADA
%----------------------------------
\begin{frame}
    \titlepage
    \note{
        Bienvenida a los participantes. Introducir el objetivo de la sesión: aprender los fundamentos de la PRL y su relevancia en el ámbito laboral, considerando también los aspectos medioambientales.
    }
\end{frame}

%----------------------------------
% ÍNDICE
%----------------------------------
\begin{frame}{Agenda}
    \tableofcontents
    \note{
        Explicar brevemente los temas que se cubrirán, destacando la importancia de cada uno para entender la PRL y la protección medioambiental.
    }
\end{frame}

%%%%%%%%%%%%%%%%%%%%%%%%%%%%%%%%%%%%%%%%%
\section{Introducción a la PRL}
%%%%%%%%%%%%%%%%%%%%%%%%%%%%%%%%%%%%%%%%%

%----------------------------------
\begin{frame}{¿Qué es la Prevención de Riesgos Laborales (PRL)?}
    \begin{itemize}
        \item Conjunto de medidas destinadas a evitar riesgos para la salud y la seguridad en el entorno laboral.
        \item Prevención vs. reacción: una inversión en seguridad que también evita consecuencias legales y económicas.
        \item Rol clave en todos los sectores y su relación con la protección del medioambiente.
    \end{itemize}
    \note{
        Ejemplo: Comparar la PRL con un seguro: invertir en seguridad resulta menos costoso que afrontar las consecuencias de un accidente.
    }
\end{frame}

%----------------------------------
\begin{frame}{Definiciones clave}
    \begin{itemize}
        \item \textbf{Peligro:} Es una fuente potencial de daño o una situación con capacidad para causar un efecto adverso, como lesiones, enfermedades o daños a la propiedad, al medio ambiente o a las personas.
        \item \textbf{Riesgo:} Es la probabilidad de que el peligro se materialice y cause un daño, junto con la gravedad del posible daño.
        \item \textbf{Accidente:} Evento que genera un daño.
        \item \textbf{Enfermedad laboral:} Patología derivada de la actividad profesional.
    \end{itemize}
    \\
    \vspace{1cm}
    \\
    El \textbf{objetivo} es reducir el riesgo mediante la identificación de peligros y la implementación de medidas preventivas.
    \note{
        Ejemplo: Diferenciar riesgo y peligro con una cafetera caliente: el peligro es la cafetera, el riesgo es quemarse si la tocas descuidadamente.

        Ejemplo: Un líquido inflamable tiene el peligro de arder, una máquina con partes móviles expuestas tiene el peligro de causar cortes. El riesgo asociado a un líquido inflamable dependerá de la cantidad almacenada, las condiciones de ventilación y la presencia de fuentes de ignición.

        Ejemplo: Chernobil. Es peligroso el uranio pero el riesgo aumenta a lo bestia cuando se fuerza la maquinaria para producir de más.
    }
\end{frame}



%%%%%%%%%%%%%%%%%%%%%%%%%%%%%%%%%%%%%%%%%
\section{Normativa y Legislación básica}
%%%%%%%%%%%%%%%%%%%%%%%%%%%%%%%%%%%%%%%%%

%----------------------------------
\begin{frame}{Marco Legislativo}
    \textbf{Unión Europea (UE)} 
    \begin{itemize}
        \item Tratado constitutivo de la UE: promover la mejora del medio de trabajo para proteger la salud y la seguridad de los trabajadores (Artículo 153).
        \item Directivas comunitarias integradas en la legislación de cada Estado miembro.
    \end{itemize}
    \vspace{2mm}
    \textbf{España}
    \begin{itemize}
        \item \textbf{Constitución Española (1978):} Art. 40 y 43 garantizan la seguridad y salud en el trabajo.
        \item \textbf{Estatuto de los Trabajadores (R.D. Leg. 1/1995):} Regula derechos y deberes en PRL.
        \item \textbf{Ley 31/1995 de Prevención de Riesgos Laborales:} Marco legal principal.
        \item \textbf{R.D. 39/1997 (Reglamento de los Servicios de Prevención):} Detalla obligaciones de empresarios y trabajadores.
    \end{itemize}
    \note{
        Destacar que incumplir la normativa conlleva sanciones económicas, penales y reputacionales.

     Tratado Constitutivo de la UE
El Tratado de Funcionamiento de la Unión Europea (TFUE) establece en su Artículo 153 que la UE apoyará y complementará las actividades de los Estados miembros en la mejora del entorno laboral para proteger la salud y la seguridad de los trabajadores. 


Directivas Comunitarias
La Directiva Marco 89/391/CEE es fundamental en materia de seguridad y salud en el trabajo. Esta directiva introduce medidas para promover la mejora de la seguridad y la salud de los trabajadores en el trabajo, estableciendo obligaciones para los empleadores y los trabajadores con el fin de limitar los accidentes laborales y las enfermedades profesionales. 


Además de la Directiva Marco, existen otras directivas específicas que abordan distintos aspectos de la seguridad y salud laboral, como:

Directiva 89/654/CEE: relativa a las disposiciones mínimas de seguridad y de salud en los lugares de trabajo.

Directiva 92/58/CEE: relativa a las disposiciones mínimas en materia de señalización de seguridad y de salud en el trabajo.
    }
\end{frame}

%----------------------------------
\begin{frame}{Obligaciones y Derechos}
    \textbf{Derechos de los trabajadores}
    \begin{itemize}
        \item Recibir información y formación.
        \item Derecho de ser consultado y participar en la evaluación de riesgos y prevención.
        \item Derecho a la vigilancia de la salud.
        \item Derecho a interrumpir la actividad en caso de riesgo grave e inminente.
    \end{itemize}
    \vspace{2mm}
    \textbf{Obligaciones de los trabajadores}
    \begin{itemize}
        \item Cumplir las medidas de prevención y utilizar correctamente los EPI.
        \item Informar de inmediato situaciones de riesgo.
        \item Cooperar con el empresario en materia de seguridad.
    \end{itemize}
    \note{
        Ejemplo: El trabajador debe llevar casco si la evaluación de riesgos lo exige; pero también puede exigir a la empresa que proporcione EPI en buen estado.
    }
\end{frame}

%----------------------------------
\begin{frame}{Obligaciones de la Empresa}
    \begin{itemize}
        \item Garantizar la seguridad y salud de todos los trabajadores a su servicio.
        \item Evaluar y planificar la actividad preventiva.
        \item Proporcionar equipos de trabajo seguros, planes de emergencia, señalización y formaciones.
        \item Garantizar la protección de los equipos especialmente sensibles.
        \item Más info: Ley 31/1995, de 8 de noviembre, de Prevención de Riesgos Laborales.
    \end{itemize}
    \note{
        Ejemplo: La empresa debe realizar revisiones periódicas de maquinaria y sistemas eléctricos para evitar accidentes.
        Simulacro de incendios...
    }
\end{frame}


\begin{frame}{Clasificación de los riesgos}
    \begin{itemize}
        \item \textbf{Riesgos físicos:} electricidad estática, temperaturas elevadas, radiaciones, ruido.
        \item \textbf{Riesgos eléctricos:} descargas eléctricas, cables defectuosos, sobrecargas.
        \item \textbf{Riesgos mecánicos:} caídas, golpes, cortes.
        \item \textbf{Riesgos ergonómicos:} posturas forzadas, manipulación de cargas, tareas repetitivas.
        \item \textbf{Riesgos químicos:} vapores tóxicos, baterías peligrosas.
        \item \textbf{Riesgos psicosociales:} estrés laboral, trabajo a turnos.
    \end{itemize}
\note{
Anecdota: En una sala de servidores mal ventilada, las temperaturas elevadas causaron una falla en varios equipos, interrumpiendo operaciones críticas por horas.
}
\end{frame}

\begin{frame}{Evaluación de la Peligrosidad}
    \begin{itemize}
        \item \textbf{Probabilidad de ocurrencia:} frecuencia con la que puede suceder el riesgo.
        \item \textbf{Severidad de las consecuencias:} impacto en caso de que ocurra.
    \end{itemize}
    \vspace{0.5cm}
    \textbf{Ejemplo:} Trabajar sin pulsera antiestática puede dañar componentes, lo que genera costos elevados de reparación.\\\vspace{1cm}
    \textbf{Métodos de Identificación de Riesgos}
    \begin{itemize}
        \item Inspecciones visuales.
        \item Revisión de protocolos.
        \item Uso de checklists de seguridad.
    \end{itemize}
\end{frame}
\note{
Usa una matriz de riesgos para ayudar a visualizar las áreas críticas. Explica cómo una sobrecarga eléctrica podría no ser frecuente, pero sus consecuencias son severas.

Checklist en los aviones.
}



%%%%%%%%%%%%%%%%%%%%%%%%%%%%%%%%%%%%%%%%%
\section{Identificación de Riesgos en el Montaje de Redes}
%%%%%%%%%%%%%%%%%%%%%%%%%%%%%%%%%%%%%%%%%

%----------------------------------
\begin{frame}{Análisis de riesgos en LAN}
    \textbf{Actividades básicas:}
    \begin{itemize}
        \item Instalación y mantenimiento del cableado de cobre y fibra óptica.
        \item Instalación de armarios de red y sus componentes.
        \item Instalación de canalizaciones para cables.
        \item Trabajos en altura, uso de herramientas eléctricas, manipulación de cargas.
    \end{itemize}
    \note{
        Con la documentación previa identificamos riesgos como: cortes, salpicaduras en ojos, descargas eléctricas, caídas, etc.
    }
\end{frame}

%----------------------------------
\begin{frame}{Riesgos en el trabajo con cableado de cobre}
    \begin{itemize}
        \item \textbf{Riesgo eléctrico:} A pesar de que el cableado de datos opere a voltaje bajo (2-3 V), puede haber fugas desde la red eléctrica (100-240 V).
        \item \textbf{Altos voltajes de rayos:} Riesgos por caídas de rayos en entornos exteriores.
        \item \textbf{Riesgo de corte o proyecciones:} Al pelar o cortar el cable.
        \item \textbf{Riesgo ergonómico:} Posturas forzadas al instalar o manipular cableado (o al pasar mucho tiempo en un equipo de oficina no ergonómico).
    \end{itemize}
    \note{
        Recomendación: Antes de manipular cableado, cortar la corriente y usar gafas para evitar fragmentos que salten.
    }
\end{frame}

%----------------------------------
\begin{frame}{Riesgos en el trabajo con fibra óptica}
    \begin{itemize}
        \item \textbf{Riesgo de corte y abrasión:} Al cortar fibra que contiene vidrio.
        \item \textbf{Riesgo para los ojos:} Pequeñas astillas de vidrio, luz láser o LED.
        \item \textbf{Riesgo para las vías respiratorias:} Polvo de vidrio.
        \item \textbf{Riesgo químico:} Al inhalar limpiadores o adhesivos.
    \end{itemize}
    \note{
        Uso de guantes, gafas de protección y mascarillas resulta esencial al trabajar con fibra óptica.
    }
\end{frame}

%%%%%%%%%%%%%%%%%%%%%%%%%%%%%%%%%%%%%%%%%
\section{Medidas de Prevención y Protección}
%%%%%%%%%%%%%%%%%%%%%%%%%%%%%%%%%%%%%%%%%

%----------------------------------
\begin{frame}{Técnicas de prevención}
    \begin{itemize}
        \item \textbf{Seguridad en el trabajo:} evitar accidentes mediante orden y limpieza.
        \item \textbf{Higiene industrial:} prevenir enfermedades derivadas del trabajo.
        \item \textbf{Ergonomía:} evitar la fatiga y lesiones musculares.
        \item \textbf{Psicosociología:} prevenir el estrés y la insatisfacción.
        \item \textbf{Medicina del trabajo:} vigilar y mantener la salud de los trabajadores.
        \item \textbf{Política social:} legislación y control para garantizar la seguridad.
    \end{itemize}
    \note{
        Un buen programa de prevención combina todas estas técnicas.
    }
\end{frame}

%----------------------------------
\begin{frame}{Técnicas de protección}
    \begin{itemize}
        \item \textbf{Protección colectiva (ej. barandillas, resguardos, señalización):} Protege a varios trabajadores a la vez.
        \item \textbf{Protección individual (EPI):} Protege de forma individual a cada trabajador. 
    \end{itemize}
    \textbf{Uso de los EPI como estrategia preventiva}:
    \begin{itemize}
        \item Casco, gafas, guantes, mascarilla, calzado de seguridad, etc.
        \item Cada EPI debe ir acompañado de formación y de un uso correcto.
    \end{itemize}
    \note{
        Los EPI son la “última barrera”: primero se debe intentar eliminar o reducir el riesgo por medidas colectivas.
    }
\end{frame}

%----------------------------------
\begin{frame}{Principales EPI en instalación de redes}
    \begin{itemize}
        \item \textbf{Cabeza:} Cascos de protección.
        \item \textbf{Ojos:} Gafas o pantallas faciales (contra partículas o radiaciones láser).
        \item \textbf{Oído:} Orejeras o tapones (ruido de taladros o máquinas).
        \item \textbf{Vías respiratorias:} Mascarillas (polvo de fibra, humos químicos).
        \item \textbf{Manos:} Guantes anticortes o antiestáticos.
        \item \textbf{Pies:} Calzado de seguridad (puntera reforzada, suela antideslizante y/o antiestático).
        \item \textbf{Cuerpo:} Batas, monos de trabajo; arnés en trabajos en altura.
    \end{itemize}
    \textbf{¿Protege un EPI de todos los riesgos?}
    \note{
        Elegir el EPI adecuado según la evaluación de riesgos y la normativa (marcado CE).
        El EPI es, en realidad, el último recurso. Para manipular cables eléctricos, si no hay corriente, no hacen falta aislantes. Ahora, si no hay otra forma que manipularlo mientras hay corriente... pues habrá que protegerse.
    }
\end{frame}

%%%%%%%%%%%%%%%%%%%%%%%%%%%%%%%%%%%%%%%%%
\section{Medidas de Seguridad en el Montaje LAN}
%%%%%%%%%%%%%%%%%%%%%%%%%%%%%%%%%%%%%%%%%

%----------------------------------
\begin{frame}{Prevención en trabajos con cableado}
    \textbf{Cobre:}
    \begin{itemize}
        \item Cortar los extremos de los hilos con gafas o pantalla para evitar proyecciones.
        \item Manos siempre secas, evitar zonas húmedas.
        \item Revisar conexiones a tierra y posibles fugas de corriente antes y después.
        \item No instalar cableado de cobre en tormentas eléctricas.
    \end{itemize}
    \vspace{1em}
    \textbf{Fibra óptica:}
    \begin{itemize}
        \item Guantes, gafas y mascarilla para polvos de vidrio.
        \item Identificar y desechar trozos de fibra en contenedores adecuados.
        \item Nunca mirar directamente el láser.
    \end{itemize}
\end{frame}

%----------------------------------
\begin{frame}{Prevención en armarios de red y canalizaciones}
    \begin{itemize}
        \item \textbf{Armarios de red:} riesgo de lesión muscular o aplastamiento (son pesados). Levantar con técnica adecuada (con las piernas, espalda recta). 
        \item Revisar la carga estática del armario y sus conexiones a tierra.
        \item Mantener ordenados los cables en bandejas y canalizaciones para evitar enredos.
        \item \textbf{Canalizaciones:} usar sierra o herramientas con resguardo, gafas y guantes de malla. 
        \item Revisar muros y suelos antes de taladrar (posibles conducciones de agua, gas, cableado eléctrico).
    \end{itemize}
    \note{
        Importante planificar la instalación para minimizar trabajos en altura y evitar posturas forzadas.
    }
\end{frame}

%----------------------------------
\begin{frame}{Trabajos en altura y uso de herramientas}
    \textbf{Escaleras de mano:}
    \begin{itemize}
        \item Verificar que estén homologadas, estables y suficientemente largas.
        \item Normas de separación entre escalera y pared (0,25 m por cada metro de altura).
        \item Señalizar la zona e impedir el acceso a personas ajenas al trabajo.
    \end{itemize}
    \vspace{1em}
    \textbf{Herramientas eléctricas (taladro, sierra, etc.):}
    \begin{itemize}
        \item Usar gafas, guantes, mascarilla y protección auditiva si procede.
        \item Revisar que no haya conducciones eléctricas o de gas en el punto de perforación.
        \item Apagar y desconectar antes de cambiar brocas o al terminar.
    \end{itemize}
\end{frame}

%%%%%%%%%%%%%%%%%%%%%%%%%%%%%%%%%%%%%%%%%
\section{Gestión Medioambiental}
%%%%%%%%%%%%%%%%%%%%%%%%%%%%%%%%%%%%%%%%%

%----------------------------------
\begin{frame}{Importancia de la gestión ambiental}
    \begin{itemize}
        \item Manejo adecuado de residuos (embalajes, cables, componentes electrónicos).
        \item Separar según tipo y enviar a reciclaje, reutilización o eliminación segura.
        \item Cumplir la legislación local, autonómica y nacional en materia de residuos (Ley 22/2011).
    \end{itemize}
    \note{
        Ejemplo: una batería mal desechada puede filtrar sustancias tóxicas al suelo. 
    }
\end{frame}

%----------------------------------
\begin{frame}{Normativa ambiental}
    \begin{itemize}
        \item \textbf{Directivas de la UE:} marcos generales para la gestión de residuos y protección del medioambiente.
        \item \textbf{Ley 22/2011 de Residuos y Suelos Contaminados (España):}
        \begin{itemize}
            \item Define residuos, productores y gestores.
            \item Regula el tratamiento, transporte y eliminación de RAEE (Residuos de Aparatos Eléctricos y Electrónicos).
        \end{itemize}
    \end{itemize}
    \note{
        CCAA y Entidades locales pueden exigir trámites específicos para la gestión de residuos.
    }
\end{frame}

%----------------------------------
\begin{frame}{Buenas prácticas medioambientales}
    \begin{itemize}
        \item Reducir, reutilizar y reciclar: “3R” como base de la gestión de materiales.
        \item Utilizar contenedores diferenciados para RAEE, tóner, vidrio, plástico, etc.
        \item Evitar verter líquidos o sustancias peligrosas por el desagüe.
        \item Disponer de planes de emergencia en caso de vertidos.
        \item Mantener ordenados los elementos de trabajo para prevenir accidentes ambientales (derrames, fugas, etc.).
    \end{itemize}
    \note{
        Un entorno limpio y ordenado también mejora la seguridad laboral y la imagen de la empresa.
    }
\end{frame}


\begin{frame}{Orden y Limpieza}

\hspace{1cm}\textbf{Beneficios}:
        \begin{itemize}
            \item Reducción de caídas por cables sueltos.
            \item Localización rápida de herramientas.
            \item ...
        \end{itemize}
    \\\vspace{1cm}
    \hspace{1cm}\textbf{Principios de las 5S}:
        \begin{enumerate}
            \item Seiri.
            \item Seiton.
            \item Seiso.
            \item Seiketsu.
            \item Shitsuke.
        \end{enumerate}
    \note{
        Ejemplo práctico: Implementar las 5S en un taller de reparación puede reducir un 30\% el tiempo de búsqueda de herramientas.

        Clasificar (Seiri).
        Ordenar (Seiton).
        Limpiar (Seiso).
        Estandarizar (Seiketsu).
        Mantener (Shitsuke).
        https://biblat.unam.mx/hevila/ResnonverbaGuayaquil/2017/vol7/no1/10.pdf
    }
\end{frame}

\begin{frame}
\frametitle{El origen de las 5S}
\begin{itemize}
    \item La metodología \textbf{5S} es parte del \textbf{Sistema de Producción de Toyota (TPS)}.
    \item Nació en Japón como una herramienta clave para optimizar la eficiencia y calidad en los procesos de manufactura.
    \item Su objetivo principal es mantener un entorno de trabajo organizado, seguro y eficiente.
    \item Las \textbf{5S} son fundamentales en el modelo de producción \textbf{lean}, ampliamente adoptado en empresas de todo el mundo. Sin embargo, son una fuente de inspiración para cualquier otro entorno productivo.
\end{itemize}

\end{frame}

\section{Clasificar (Seiri)}
\begin{frame}
\frametitle{1. Clasificar (Seiri)}
\begin{itemize}
    \item \textbf{Objetivo:} Identificar y separar los elementos necesarios de los innecesarios.
    \item \textbf{Qué implica:}
        \begin{itemize}
            \item Revisar el material y eliminar lo innecesario.
            \item Priorizar lo esencial para las operaciones.
        \end{itemize}
    \item \textbf{Beneficio:} Reduce el desorden y facilita el acceso a lo necesario.
    \item \textbf{Ejemplo:} Clasificar herramientas en una planta de Toyota, eliminando duplicados o elementos dañados.
\end{itemize}
\end{frame}

\section{Ordenar (Seiton)}
\begin{frame}
\frametitle{2. Ordenar (Seiton)}
\begin{itemize}
    \item \textbf{Objetivo:} Organizar los elementos necesarios en el lugar adecuado.
    \item \textbf{Qué implica:}
        \begin{itemize}
            \item Designar un lugar específico para cada cosa.
            \item Utilizar herramientas como etiquetas o estanterías.
        \end{itemize}
    \item \textbf{Beneficio:} Minimiza el tiempo perdido buscando elementos y mejora la eficiencia.
    \item \textbf{Ejemplo:} Diseñar paneles visuales para herramientas en las líneas de producción de Toyota.
\end{itemize}
\end{frame}

\section{Limpiar (Seiso)}
\begin{frame}
\frametitle{3. Limpiar (Seiso)}
\begin{itemize}
    \item \textbf{Objetivo:} Mantener el lugar de trabajo limpio y en condiciones óptimas.
    \item \textbf{Qué implica:}
        \begin{itemize}
            \item Realizar limpieza regular de equipos y herramientas.
            \item Identificar y corregir fuentes de suciedad.
        \end{itemize}
    \item \textbf{Beneficio:} Mejora la seguridad y el ambiente de trabajo.
    \item \textbf{Ejemplo:} Inspeccionar y limpiar vehículos antes de entregarlos en Toyota.
\end{itemize}
\end{frame}

\section{Estandarizar (Seiketsu) y Mantener (Shitsuke)}
\begin{frame}
\frametitle{4. Estandarizar (Seiketsu) y Mantener (Shitsuke)}
\begin{itemize}
    \item \textbf{Estandarizar (Seiketsu):}
        \begin{itemize}
            \item Crear normas y protocolos claros para mantener las tres primeras S.
            \item Usar listas de verificación y señalización visual.
        \end{itemize}
    \item \textbf{Mantener (Shitsuke):}
        \begin{itemize}
            \item Fomentar hábitos y promover la autorregulación.
            \item Realizar auditorías periódicas y capacitaciones.
        \end{itemize}
    \item \textbf{Beneficio:} Asegura que las mejoras sean sostenibles y parte de la cultura organizacional.
    \item \textbf{Ejemplo:} En Toyota, auditorías regulares aseguran la disciplina en el uso de las 5S.
\end{itemize}
\end{frame}

%%%%%%%%%%%%%%%%%%%%%%%%%%%%%%%%%%%%%%%%%
\section{Conclusiones}
%%%%%%%%%%%%%%%%%%%%%%%%%%%%%%%%%%%%%%%%%

\begin{frame}{Conclusiones}
    \begin{itemize}
        \item La PRL es un conjunto de técnicas y obligaciones legales que busca proteger la salud de los trabajadores.
        \item En el montaje y mantenimiento de redes LAN, existen riesgos específicos (eléctricos, mecánicos, químicos, ergonómicos, etc.) que debemos prevenir.
        \item El uso de EPI, junto a la protección colectiva y la formación, reduce significativamente los accidentes.
        \item La gestión de residuos y la protección medioambiental forman parte integral de la prevención en el trabajo.
        \item Cumplir con la normativa evita sanciones y contribuye a un entorno laboral y ambiental más seguro.
    \end{itemize}
    \note{
        Destacar que la PRL y el cuidado ambiental son inversiones a largo plazo: aumentan la productividad y la reputación de la organización.
    }
\end{frame}

%----------------------------------
\begin{frame}
    \centering
    \Huge ¿Preguntas?
\end{frame}

\end{document}
